\section{Statement and setup}

We restate the proposition to be proved and then give a complete, rigorous proof.

\begin{introtheorem}
There is a constant $k>0$ so that the following holds. Suppose that $R\subset\RR^4$ is a $4$-dimensional rectangle with side lengths $R_1\le R_2\le R_3\le R_4$ and $S$ is a $4$-dimensional rectangle with side lengths $S_1\le S_2\le S_3\le S_4$. Suppose that $f:(R,\partial R)\to(S,\partial S)$ is a piecewise smooth map of degree $1$ with $\Dil_2(f)\le 1$. Suppose that $R_1\le kS_1$. Then either
\[
R_3R_4>kS_3S_4,\qquad\text{or}\qquad R_1R_2R_3R_4>kS_1S_2^{1/2}S_3^{3/2}S_4 .
\]
\end{introtheorem}

Throughout, $\Vol_k$ denotes $k$-dimensional volume (mass), and $\Dil_k(f)$ is the $k$-dilation of $f$, equivalently $\sup_{x}\|\Lambda^k df_x\|$. We write $A\lesssim B$ if $A\le CB$ for an absolute constant $C$, and $A\sim B$ if $A\lesssim B$ and $B\lesssim A$. The constants implicit in $\lesssim$ are absolute (in particular independent of all $R_i$, $S_i$).

We will use two ingredients: an isoperimetric inequality for relative cycles in a rectangle, and a homotopy-invariant linking number. We develop both, then combine them.

\section{The isoperimetric lemma}

Let $R=[0,R_1]\times\cdots\times[0,R_n]$ with $R_1\le\cdots\le R_n$. A \emph{relative integral $k$-cycle} $z$ in $R$ is an integral Lipschitz $k$-chain in $R$ with $\partial z$ supported in $\partial R$. The \emph{filling volume} of $z$ is the infimal mass of a relative $(k+1)$-chain $y$ in $R$ with $\partial y - z$ supported in $\partial R$. Let $I_R^k(V)$ denote the supremum of the filling volume over all relative integral $k$-cycles $z$ in $R$ with $\Vol_k(z)\le V$.

\begin{lemma}\label{iso0}
$I_R^0(V)\le R_1V$.
\end{lemma}

\begin{proof}
A relative $0$-cycle is a finite integer combination $\sum_i c_i p(i)$ of points $p(i)\in R$ with $\Vol_0(z)=\sum_i|c_i|$. A point $p=(p_1,\dots,p_n)$ bounds (modulo $\partial R$) the segment $[0,p_1]\times\{p_2\}\times\cdots\times\{p_n\}$, which has length $p_1\le R_1$ and whose $0$-boundary is $p$ together with a point on $\partial R$. Taking this segment with multiplicity $c_i$ for each $i$ yields a relative $1$-chain $y$ with $\partial y \equiv z$ modulo $\partial R$ and $\Vol_1(y)\le R_1\sum_i|c_i| = R_1\Vol_0(z)$.
\end{proof}

We recall the Federer--Fleming deformation theorem in the form we need; a proof appears in \cite{G3}.

\begin{theorem}[Federer--Fleming deformation]\label{FF}
Let $z$ be a relative integral $k$-cycle in $R$. Fix $L\le R_1$ and a rectangular lattice in $R$ whose cells have side lengths between $L$ and $2L$ in each coordinate, with $\partial R$ contained in the $(n-1)$-skeleton of the lattice. Then there is a relative integral $k$-cycle $z'$ contained in the $k$-skeleton of the lattice with
\begin{enumerate}
\item $\Vol_k(z')\le C(n)\,\Vol_k(z)$;
\item the filling volume of $z'-z$ in $R$ is at most $C(n)\,L\,\Vol_k(z)$.
\end{enumerate}
Moreover, the lattice in the $x_1$ direction may be chosen with cell boundaries only at $x_1=0$ and $x_1=R_1$ when $L=R_1$.
\end{theorem}

\begin{lemma}[Isoperimetric lemma]\label{iso}
There are constants $c(n)>0$, $C(n)$ so that for every $1\le k\le n-1$ the following holds. If $V\le c(n)R_1\cdots R_k$, write
\[
V=c(n)R_1\cdots R_j\,\rho^{\,k-j}
\]
for the unique $0\le j\le k-1$ and $\rho$ with $R_j\le\rho\le R_{j+1}$ (where $R_0:=0$, so $j=0$ forces $\rho\le R_1$ and the prefactor is $c(n)$). Then
\[
I_R^k(V)\le C(n)\,R_1\cdots R_j\,\rho^{\,k-j+1}.
\]
In every case (regardless of the size of $V$), $I_R^k(V)\le C(n)\,R_{k+1}\,V$.
\end{lemma}

\begin{proof}
We induct on $k$, with Lemma \ref{iso0} as the base case $k=0$ in the inductive step below (and the bound $I^0_R(V)\le R_1V$ already proved). Fix $k\ge 1$ and assume the statement for $k-1$ in every rectangle of dimension $n-1$. Let $z$ be a relative integral $k$-cycle in $R$ with $\Vol_k(z)=V$.

\smallskip\noindent\emph{Case 1: $V\le c(n)R_1^k$.}
Choose $L=C(n)V^{1/k}$. For $c(n)$ small and $C(n)$ chosen suitably, $L\le R_1$. Apply Theorem \ref{FF} with this $L$ to get $z'$ in the $k$-skeleton with $\Vol_k(z')\le C(n)V$. Since each $k$-cell of the lattice has volume at least $L^k\ge C(n)V$ (with $C(n)$ large enough), and $z'$ is an integer combination of such cells with total mass $\le C(n)V<L^k$, every coefficient must vanish: $z'=0$. By Theorem \ref{FF}(2), $z=z'+(z-z')$ is filled by a chain of mass at most $C(n)LV\le C(n)V^{(k+1)/k}$. This is exactly the claimed bound when $V\le c(n)R_1^k$, i.e.\ when $j=0$ and $\rho=V^{1/k}/c(n)^{1/k}\sim V^{1/k}\le R_1$: indeed $R_1\cdots R_j\rho^{k-j+1}=\rho^{k+1}\sim V^{(k+1)/k}$.

\smallskip\noindent\emph{Case 2: $V\ge c(n)R_1^k$.}
Choose $L=R_1$ and a lattice as in Theorem \ref{FF} with cell boundaries in the $x_1$-direction only at $x_1=0,R_1$. We obtain $z'$ in the $k$-skeleton with $\Vol_k(z')\le C(n)V$ and a filling of $z-z'$ of mass at most $C(n)R_1V$. Each interior $k$-cell of the lattice that can carry nonzero coefficient is of the form $[0,R_1]\times(\text{a }(k-1)\text{-cell of the lattice on }[0,R_2]\times\cdots\times[0,R_n])$, because in the $x_1$ direction the only lattice interval is $[0,R_1]$. (Cells lying in $\partial R$ are irrelevant modulo $\partial R$.) Hence
\[
z'=[0,R_1]\times z_1,
\]
where $z_1$ is a relative integral $(k-1)$-cycle in the $(n-1)$-rectangle $R':=[0,R_2]\times\cdots\times[0,R_n]$ with
\[
\Vol_{k-1}(z_1)=\frac{\Vol_k(z')}{R_1}\le \frac{C(n)V}{R_1}.
\]

\emph{Filling $z_1$.} If $k=1$, fill $z_1$ (a relative $0$-cycle in $R'$) by Lemma \ref{iso0}; if $k\ge2$, fill $z_1$ by the inductive hypothesis in $R'$. Call the filling $C_1$, a relative $k$-chain in $R'$. Then $[0,R_1]\times C_1$ is a relative $(k+1)$-chain in $R$ with
\[
\partial\big([0,R_1]\times C_1\big)\equiv [0,R_1]\times z_1=z' \pmod{\partial R},
\]
and $\Vol_{k+1}\big([0,R_1]\times C_1\big)=R_1\,\Vol_k(C_1)$. Adding the filling of $z-z'$ of mass $\le C(n)R_1V$ gives a filling of $z$.

\emph{Bookkeeping.} First the universal bound. By induction (or Lemma \ref{iso0} for $k=1$), $\Vol_k(C_1)\le C(n)R_{k+1}\,\Vol_{k-1}(z_1)\le C(n)R_{k+1}V/R_1$ (the second factor of the rectangle $R'$ has side lengths $R_2\le\cdots\le R_n$, so its ``$R'_{k}$'' is $R_{k+1}$). Therefore
\[
\Vol_{k+1}\big([0,R_1]\times C_1\big)\le C(n)R_{k+1}V,
\]
and adding $C(n)R_1V\le C(n)R_{k+1}V$ gives $I_R^k(V)\le C(n)R_{k+1}V$ in all cases.

Now the refined bound when $V\le c(n)R_1\cdots R_k$ (and $V\ge c(n)R_1^k$, so $j\ge1$). Write $V=c(n)R_1\cdots R_j\rho^{k-j}$ with $R_j\le\rho\le R_{j+1}$ and $1\le j\le k-1$. Then
\[
\Vol_{k-1}(z_1)\le \frac{C(n)V}{R_1}= C(n)\,c(n)\,R_2\cdots R_j\,\rho^{(k-1)-(j-1)} .
\]
Thus $z_1$ is a relative $(k-1)$-cycle in $R'$ whose volume is, up to the constant, of the form $c(n-1)R_2\cdots R_j\,\rho^{(k-1)-(j-1)}$ with $R_j\le\rho\le R_{j+1}$; the same $j-1$ and $\rho$ are the data of Lemma \ref{iso} applied to $R'$ and dimension $k-1$ (after adjusting $c(n-1),C(n-1)$ to absorb the constant factor). By the inductive hypothesis,
\[
\Vol_k(C_1)\le C(n-1)\,R_2\cdots R_j\,\rho^{(k-1)-(j-1)+1}=C(n-1)\,R_2\cdots R_j\,\rho^{\,k-j+1}.
\]
Hence
\[
\Vol_{k+1}\big([0,R_1]\times C_1\big)\le C(n)\,R_1\cdots R_j\,\rho^{\,k-j+1}.
\]
Finally, the filling of $z-z'$ has mass $\le C(n)R_1V=C(n)\,R_1\cdot c(n)R_1\cdots R_j\rho^{k-j}\le C(n)R_1\cdots R_j\rho^{k-j+1}$, using $R_1\le R_j\le\rho$. Adding the two contributions yields
\[
I_R^k(V)\le C(n)\,R_1\cdots R_j\,\rho^{\,k-j+1},
\]
as claimed. This completes the induction.
\end{proof}

We record the two special cases we will use, in $n=4$, $k=2$.

\begin{lemma}\label{iso2}
There are constants $c>0$, $C$ so that the following holds for the $4$-rectangle $R$.
\begin{enumerate}
\item If $y$ is a relative integral $2$-cycle in $R$ with $cR_1^2\le \Vol_2(y)\le cR_1R_2$, then $y$ bounds a relative $3$-chain $z$ in $R$ with
\[
\Vol_3(z)\le C\,R_1^{-1}\Vol_2(y)^2 .
\]
\item If $y$ is a relative integral $2$-cycle in $R$ with $\Vol_2(y)\le cR_1^2$, then $y$ bounds a relative $3$-chain $z$ in $R$ with
\[
\Vol_3(z)\le C\,\Vol_2(y)^{3/2}\le C\,R_1^3 .
\]
\end{enumerate}
\end{lemma}

\begin{proof}
Apply Lemma \ref{iso} with $n=4$, $k=2$, $V=\Vol_2(y)$.

(2) If $V\le cR_1^2$ then $j=0$, $\rho\sim V^{1/2}\le R_1$, and the bound is $I_R^2(V)\le C\rho^3\sim CV^{3/2}\le CR_1^3$.

(1) If $cR_1^2\le V\le cR_1R_2$ then $j=1$ and $V=c\,R_1\rho$ with $R_1\le\rho\le R_2$, i.e.\ $\rho\sim V/R_1$. The bound is
\[
I_R^2(V)\le C\,R_1\,\rho^{\,2}\sim C\,R_1\Big(\frac{V}{R_1}\Big)^2=C\,R_1^{-1}V^2 .\qedhere
\]
\end{proof}

\section{The linking invariant}

Let $*$ be the wedge point of $S^2\vee S^2$, and $\Omega_1,\Omega_2\in H^2(S^2\vee S^2;\ZZ)$ the generators dual to the two spheres. Choose closed differential forms $\omega_1,\omega_2$ representing $\Omega_1,\Omega_2$, with $\omega_i$ supported in the $i$-th sphere away from a neighborhood of $*$. Then $\omega_1\wedge\omega_2=0$ identically (their supports are disjoint).

Let $(M,\partial M)$ be a compact piecewise smooth manifold with boundary, $h\in H_3(M,\partial M;\ZZ)$, and let $g:(M,\partial M)\to(S^2\vee S^2,*)$ be piecewise smooth with $g^*\Omega_i=0$ in $H^2(M,\partial M;\ZZ)$ for $i=1,2$. By the de Rham theorem for relative (compactly supported) cohomology, there exist piecewise smooth $1$-forms $\alpha_1,\alpha_2$ on $M$, each supported in the interior, with $d\alpha_i=g^*\omega_i$. For a relative integral $3$-chain $z$ with $[z]=h$ we set
\[
L_h(g):=\int_z \alpha_1\wedge g^*\omega_2 .
\]
If $M$ is a closed-up oriented $3$-manifold with fundamental class $[M]$ we write $L(g):=L_{[M]}(g)$.

\begin{lemma}\label{link}
Assume $g^*\Omega_i=0$ in $H^2(M,\partial M;\ZZ)$ for $i=1,2$. Then:
\begin{enumerate}
\item $L_h(g)$ is independent of the choice of $z$ and of the $\alpha_i$.
\item If $\beta_1$ is any $1$-form on $M$ (whose support may meet $\partial M$) with $d\beta_1=g^*\omega_1$, then $L_h(g)=\int_z\beta_1\wedge g^*\omega_2$.
\item $L_h(g)$ is invariant under homotopies of $g$ through maps of pairs.
\item If $\phi:(L,\partial L)\to(M,\partial M)$ with $\phi_*h_L=h_M$, then $L_{h_L}(g\circ\phi)=L_{h_M}(g)$.
\end{enumerate}
\end{lemma}

\begin{proof}
The $3$-form $\alpha_1\wedge g^*\omega_2$ is closed: $d(\alpha_1\wedge g^*\omega_2)=g^*\omega_1\wedge g^*\omega_2=g^*(\omega_1\wedge\omega_2)=0$.

(1) Independence of $z$: if $[z_1]=[z_2]=h$ then $z_1-z_2=\partial y+w$ with $w$ supported in $\partial M$ and $y$ a relative $4$-chain. Since $\alpha_1\wedge g^*\omega_2$ is closed and $\alpha_1$ (hence the integrand) is supported in the interior, Stokes' theorem gives $\int_{z_1}=\int_{z_2}$ (the contribution of $w\subset\partial M$ vanishes because $\alpha_1$ vanishes near $\partial M$). Independence of $\alpha_2$: if $d\alpha_2=d\alpha_2'=g^*\omega_2$ with both interior-supported, then $\alpha_2-\alpha_2'$ is closed and interior-supported, and $\int_z\alpha_1\wedge d(\alpha_2-\alpha_2')$ etc. is handled by (2) below; more directly, independence of $\alpha_1$ is (2).

(2) Let $d\beta_1=d\alpha_1=g^*\omega_1$. Then
\[
\int_z(\beta_1-\alpha_1)\wedge g^*\omega_2=\int_z(\beta_1-\alpha_1)\wedge d\alpha_2=\int_z d\big((\beta_1-\alpha_1)\wedge\alpha_2\big),
\]
using $d(\beta_1-\alpha_1)=0$. By Stokes, this equals the integral over $\partial z$, which lies in $\partial M$; there $\alpha_2$ vanishes (it is interior-supported), so the integral is $0$. Hence $\int_z\beta_1\wedge g^*\omega_2=\int_z\alpha_1\wedge g^*\omega_2=L_h(g)$. Taking $\beta_1=\alpha_1'$ another interior primitive shows independence of $\alpha_1$ in (1).

(3) Let $G:(M\times[0,1],\partial M\times[0,1])\to(S^2\vee S^2,*)$ be a homotopy from $g_0$ to $g_1$. Pick interior-supported $1$-forms $A_i$ on $M\times[0,1]$ with $dA_i=G^*\omega_i$. As in (1), $A_1\wedge G^*\omega_2$ is closed and interior-supported. Stokes' theorem on $Z=z\times[0,1]$ (whose boundary is $z\times\{1\}-z\times\{0\}$ plus a piece in $\partial M\times[0,1]$ where $A_1$ vanishes) gives
\[
0=\int_{\partial Z}A_1\wedge G^*\omega_2=\int_{z\times\{1\}}-\int_{z\times\{0\}}=L_h(g_1)-L_h(g_0).
\]

(4) Choose interior-supported $\alpha_{i,M}$ on $M$ with $d\alpha_{i,M}=g^*\omega_i$. Then $\alpha_{i,L}:=\phi^*\alpha_{i,M}$ is interior-supported on $L$ with $d\alpha_{i,L}=(g\circ\phi)^*\omega_i$. Choose $z_L$ with $[z_L]=h_L$ and let $z_M=\phi_\#z_L$, so $[z_M]=\phi_*h_L=h_M$. Then
\[
L_{h_L}(g\circ\phi)=\int_{z_L}\phi^*\big(\alpha_{1,M}\wedge g^*\omega_2\big)=\int_{z_M}\alpha_{1,M}\wedge g^*\omega_2=L_{h_M}(g).\qedhere
\]
\end{proof}

\paragraph{A model map with $L=1$.}
Let $Q_1=[0,1]^3$. There is a piecewise smooth map $g_1:(Q_1,\partial Q_1)\to(S^2\vee S^2,*)$ with $L(g_1)=1$ (a standard construction; see \cite{GH}, Prop.\ 5.2). Equip $S^2\vee S^2$ with the metric making each sphere round of radius $S_1$, and choose $\omega_i$ representing $\Omega_i$ with
\[
\|\omega_i\|_{L^\infty}\lesssim S_1^{-2}.
\]
Let $Q_{S_1}$ be a cube of side length $S_1$. Rescaling $g_1$ by the dilation $Q_{S_1}\to Q_1$ produces a map
\[
g_{S_1}:(Q_{S_1},\partial Q_{S_1})\to(S^2\vee S^2,*),\qquad \Dil_1(g_{S_1})\lesssim 1,\quad L(g_{S_1})=1 .
\]
(The linking number $L$ is a homotopy invariant and is unchanged under reparametrization of the domain; the $1$-dilation bound holds because each round sphere of radius $S_1$ is reached by a map of $1$-dilation $\lesssim 1$ from a cube of side $S_1$.)

\section{Proof of the Proposition}

Let $f:(R,\partial R)\to(S,\partial S)$ have degree $1$ and $\Dil_2(f)\le1$, with $R_1\le kS_1$ for $k$ to be fixed.

\subsection*{Slices of the target and a good set of parameters}
For $(a,b)\in[0,S_3]\times[0,S_4]$ let
\[
\Sigma_{a,b}:=[0,S_1]\times[0,S_2]\times\{a\}\times\{b\}\subset S .
\]
By the transversality (Sard) theorem, for almost every $(a,b)$ the map $f$ is transverse to $\Sigma_{a,b}$, so $f^{-1}(\Sigma_{a,b})$ is a piecewise smooth $2$-submanifold with boundary in $\partial R$. Define
\[
A:=\Big\{(a,b)\in[\tfrac14S_3,\tfrac34S_3]\times[\tfrac14S_4,\tfrac34S_4]:\ f\pitchfork\Sigma_{a,b}\Big\}.
\]
Then $|A|=\tfrac14 S_3\cdot\tfrac14 S_4=\tfrac1{16}S_3S_4\sim S_3S_4$, since the transversality condition fails only on a null set.

\subsection*{Topology of $(S,\partial S\cup\Sigma_{a,b})$}
Fix $(a,b)\in A$. Since $\partial\Sigma_{a,b}\subset\partial S$ is null-homologous in $\partial S$, choose a $2$-chain $y\subset\partial S$ with $\partial y=-\partial\Sigma_{a,b}$ and set $z_{2,S}=y+\Sigma_{a,b}$, a $2$-cycle in $\partial S\cup\Sigma_{a,b}$. One checks $H_2(\partial S\cup\Sigma_{a,b};\ZZ)\cong\ZZ$, generated by $h_{2,S}:=[z_{2,S}]$. Likewise $H_3(S,\partial S\cup\Sigma_{a,b};\ZZ)\cong\ZZ$: choosing a relative $3$-chain $z_{3,S}$ in $(S,\partial S)$ with $\partial z_{3,S}=\Sigma_{a,b}$, the class $h_{3,S}:=[z_{3,S}]$ generates it, and the connecting map $\partial:H_3(S,\partial S\cup\Sigma_{a,b})\to H_2(\partial S\cup\Sigma_{a,b})$ sends $h_{3,S}\mapsto h_{2,S}$ and is an isomorphism (both groups are $\ZZ$ and this map is onto by construction).

\subsection*{Construction of a high-linking map on $S$}
Let $S'=[0,S_1]\times[0,S_2]\times[0,S_3]$. Define $r:[0,S_3]\times[0,S_4]\to[0,S_3]$ by
\[
r(y_3,y_4)=\min\big(10(|y_3-a|+|y_4-b|),\,S_3\big),
\]
and $g_1(y_1,y_2,y_3,y_4)=(y_1,y_2,r(y_3,y_4))$, so $g_1:S\to S'$. Then $g_1(\Sigma_{a,b})\subset[0,S_1]\times[0,S_2]\times\{0\}\subset\partial S'$, and $g_1(\partial S)\subset\partial S'$ (the faces $y_1\in\{0,S_1\}$ or $y_2\in\{0,S_2\}$ map to corresponding faces; the faces $y_3,y_4\in\{0,S_4\text{ etc.}\}$ map into the face $\{r=S_3\}$ when $(a,b)$ is interior, which holds since $(a,b)\in A$). Thus
\[
g_1:(S,\partial S\cup\Sigma_{a,b})\to(S',\partial S')
\]
is a map of pairs. Moreover $g_{1,*}(h_{3,S})=[S']$, the fundamental class of $(S',\partial S')$: indeed $g_1$ restricted to $z_{3,S}$ is degree $1$ onto $S'$ because $r$ achieves the value $0$ transversally exactly at $(a,b)$, an interior point.

Choose disjoint cubes $Q_1,\dots,Q_N\subset S'$ of side length $S_1$ with $N\sim S_2S_3/S_1^2$ (possible since $S'$ has dimensions $S_1\le S_2\le S_3$, giving $\lfloor S_2/S_1\rfloor\lfloor S_3/S_1\rfloor\sim S_2S_3/S_1^2$ such cubes stacked in the $y_2,y_3$ directions; here we use $S_1\le S_2\le S_3$). On each $Q$ use the model map $g_Q:=g_{S_1}$ (translated), with $L(g_Q)=1$, $\Dil_1(g_Q)\lesssim 1$, and extend by the constant $*$ outside $\bigcup_\ell Q_\ell$ to get
\[
g_2:(S',\partial S')\to(S^2\vee S^2,*),\qquad \Dil_1(g_2)\lesssim 1 .
\]
Since $H^2(S',\partial S';\ZZ)=0$, $g_2^*\Omega_i=0$ and $L(g_2)$ is defined. Choosing interior primitives $\alpha_{1,Q}$ on each $Q$ with $d\alpha_{1,Q}=g_Q^*\omega_1$ and gluing to $\alpha_{1,2}$ (zero off $\bigcup Q_\ell$), we compute
\[
L(g_2)=\int_{S'}\alpha_{1,2}\wedge g_2^*\omega_2=\sum_\ell\int_{Q_\ell}\alpha_{1,Q_\ell}\wedge g_{Q_\ell}^*\omega_2=\sum_\ell L(g_{Q_\ell})=N\sim\frac{S_2S_3}{S_1^2}.
\]
Set $g:=g_2\circ g_1:(S,\partial S\cup\Sigma_{a,b})\to(S^2\vee S^2,*)$. Then $\Dil_1(g)\lesssim1$, hence $\Dil_2(g)\lesssim1$, and by Lemma \ref{link}(4) (with $\phi=g_1$, $h_{3,S}\mapsto[S']$),
\[
L_{h_{3,S}}(g)=L(g_2)\sim\frac{S_2S_3}{S_1^2}.
\]

\subsection*{Pulling back through $f$}
View $f$ as a map of pairs $f:(R,\partial R\cup f^{-1}(\Sigma_{a,b}))\to(S,\partial S\cup\Sigma_{a,b})$. As above define
\[
h_{2,R}\in H_2(R,\partial R\cup f^{-1}(\Sigma_{a,b});\ZZ)
\]
by $z_{2,R}=y_R+f^{-1}(\Sigma_{a,b})$ with $y_R\subset\partial R$, $\partial y_R=-\partial f^{-1}(\Sigma_{a,b})$. Because $f:(R,\partial R)\to(S,\partial S)$ has degree $1$ and $f\pitchfork\Sigma_{a,b}$, the restriction $f:(f^{-1}(\Sigma_{a,b}),\partial(\cdot))\to(\Sigma_{a,b},\partial\Sigma_{a,b})$ has degree $1$, so $f_*h_{2,R}=h_{2,S}$.

Choose a relative $3$-chain $z_3$ in $(R,\partial R)$ with $\partial z_3=f^{-1}(\Sigma_{a,b})$ (such $z_3$ exists since $[f^{-1}(\Sigma_{a,b})]=0$ in $H_2(R,\partial R)$, because its image $\Sigma_{a,b}$ bounds $z_{3,S}$ and $f$ has degree $1$; equivalently $H_2(R,\partial R)=0$). View $z_3$ as a relative $3$-cycle in $(R,\partial R\cup f^{-1}(\Sigma_{a,b}))$ and set $h_{3,R}:=[z_3]$. Then $\partial h_{3,R}=h_{2,R}$, hence
\[
\partial\big(f_*h_{3,R}\big)=f_*h_{2,R}=h_{2,S}.
\]
Since $\partial:H_3(S,\partial S\cup\Sigma_{a,b})\to H_2(\partial S\cup\Sigma_{a,b})$ is an isomorphism and $\partial h_{3,S}=h_{2,S}$, we conclude $f_*h_{3,R}=h_{3,S}$.

Let $F:=g\circ f:(R,\partial R\cup f^{-1}(\Sigma_{a,b}))\to(S^2\vee S^2,*)$. Then $\Dil_2(F)\le\Dil_2(g)\,\Dil_2(f)\lesssim1$. By Lemma \ref{link}(4) with $\phi=f$,
\begin{equation}\label{eq:linkF}
L_{h_{3,R}}(F)=L_{h_{3,S}}(g)\sim\frac{S_2S_3}{S_1^2}.
\end{equation}

\subsection*{Upper bound for the linking number via $\Dil_2$}
Since $\|\omega_i\|_{L^\infty}\lesssim S_1^{-2}$ and $\Dil_2(F)\lesssim1$, the pulled-back $2$-forms satisfy
\[
\|F^*\omega_i\|_{L^\infty}\lesssim S_1^{-2}\qquad(i=1,2),
\]
because for a unit-norm simple $2$-vector $v$, $|F^*\omega_i(v)|=|\omega_i(\Lambda^2dF\,v)|\le\|\omega_i\|_\infty\|\Lambda^2dF\|\lesssim S_1^{-2}$.

Construct a primitive $\beta_1$ of $F^*\omega_1$ by integrating in $x_1$. Writing
\[
F^*\omega_1=\sum_{j=2}^4 c_j(x)\,dx_1\wedge dx_j+\sum_{2\le j_1<j_2\le4}c_{j_1j_2}(x)\,dx_{j_1}\wedge dx_{j_2},
\]
set
\[
\beta_1=\sum_{j=2}^4\Big(\int_0^{x_1}c_j(x_1',x_2,x_3,x_4)\,dx_1'\Big)dx_j .
\]
Since $F^*\omega_1$ is closed, the Poincaré-lemma computation (see \cite{BT}, Ch.\ 1) gives $d\beta_1=F^*\omega_1$. As $|c_j|\le\|F^*\omega_1\|_\infty\lesssim S_1^{-2}$ and $0\le x_1\le R_1$,
\[
\|\beta_1\|_{L^\infty}\lesssim R_1S_1^{-2}.
\]
By Lemma \ref{link}(2),
\[
|L_{h_{3,R}}(F)|=\Big|\int_{z_3}\beta_1\wedge F^*\omega_2\Big|\le \Vol_3(z_3)\,\|\beta_1\|_\infty\,\|F^*\omega_2\|_\infty\lesssim \Vol_3(z_3)\,R_1S_1^{-4}.
\]
Combining with \eqref{eq:linkF},
\[
\frac{S_2S_3}{S_1^2}\lesssim \Vol_3(z_3)\,R_1S_1^{-4}\quad\Longrightarrow\quad
\boxed{\ R_1^{-1}S_1^2S_2S_3\lesssim \Vol_3(z_3).\ }
\tag{$\ast$}
\]

\subsection*{Lower bound for $\Vol_2(f^{-1}(\Sigma_{a,b}))$}
We now bound $\Vol_3(z_3)$ from above using the isoperimetric lemma; comparing with $(\ast)$ forces $\Vol_2(f^{-1}(\Sigma_{a,b}))$ to be large. Recall $z_3$ may be taken to be a minimal filling of $y:=f^{-1}(\Sigma_{a,b})$, so $\Vol_3(z_3)\le I_R^2(\Vol_2(y))$.

First we rule out case (2) of Lemma \ref{iso2}. From $(\ast)$ and $R_1\le kS_1$ (so $S_1\ge k^{-1}R_1$),
\[
\Vol_3(z_3)\gtrsim R_1^{-1}S_1^2S_2S_3\ge R_1^{-1}S_1^2\cdot S_1\cdot S_1 = R_1^{-1}S_1^4\ge k^{-4}R_1^3,
\]
using $S_2\ge S_1$, $S_3\ge S_1$. Choosing $k>0$ small enough that $k^{-4}>C$ (the constant $C$ of Lemma \ref{iso2}(2)), the bound $\Vol_3(z_3)\le CR_1^3$ in \eqref{iso2}(2) is violated. Hence the hypothesis $\Vol_2(y)\le cR_1^2$ of \eqref{iso2}(2) must fail:
\[
\Vol_2(f^{-1}(\Sigma_{a,b}))> cR_1^2 .
\]
Consequently, by Lemma \ref{iso2}, only two possibilities remain for each $(a,b)\in A$:

\emph{(i)} $\Vol_2(f^{-1}(\Sigma_{a,b}))\ge cR_1R_2$ (the range \eqref{iso2}(1) is exceeded on the right); or

\emph{(ii)} $cR_1^2\le\Vol_2(y)\le cR_1R_2$, in which case Lemma \ref{iso2}(1) gives $\Vol_3(z_3)\le CR_1^{-1}\Vol_2(y)^2$, and combining with $(\ast)$,
\[
R_1^{-1}S_1^2S_2S_3\lesssim R_1^{-1}\Vol_2(y)^2
\quad\Longrightarrow\quad
\Vol_2(f^{-1}(\Sigma_{a,b}))\gtrsim S_1S_2^{1/2}S_3^{1/2}.
\]
In summary, for every $(a,b)\in A$,
\begin{equation}\label{eq:two}
\Vol_2(f^{-1}(\Sigma_{a,b}))\gtrsim R_1R_2,
\qquad\text{or}\qquad
\Vol_2(f^{-1}(\Sigma_{a,b}))\gtrsim S_1S_2^{1/2}S_3^{1/2}.
\end{equation}

\subsection*{Coarea and conclusion}
The slices $\{\Sigma_{a,b}\}$ foliate $S$ over the $(y_3,y_4)$-plane; their preimages $f^{-1}(\Sigma_{a,b})$ are the level sets of the projection $\pi\circ f$, where $\pi:S\to[0,S_3]\times[0,S_4]$ is $(y_1,y_2,y_3,y_4)\mapsto(y_3,y_4)$. Since $\Dil_2(f)\le1$, the map $\pi\circ f:R\to\RR^2$ has $2$-dilation $\le1$ (as $\pi$ is $1$-Lipschitz and does not increase $2$-volume of the $\{y_3,y_4\}$-block), so by the coarea inequality
\[
R_1R_2R_3R_4=\Vol_4(R)\ \ge\ \int_{[0,S_3]\times[0,S_4]}\Vol_2\big(f^{-1}(\Sigma_{a,b})\big)\,da\,db
\ \ge\ \int_A \Vol_2\big(f^{-1}(\Sigma_{a,b})\big)\,da\,db .
\]
Partition $A=A_1\cup A_2$, where $A_1$ is the set of $(a,b)\in A$ obeying the first alternative of \eqref{eq:two} and $A_2$ the second; at least one has measure $\ge\frac12|A|\sim S_3S_4$. If $|A_1|\gtrsim S_3S_4$ then
\[
R_1R_2R_3R_4\ \gtrsim\ R_1R_2\,|A_1|\ \gtrsim\ R_1R_2\,S_3S_4,
\]
which (dividing by $R_1R_2$) gives $R_3R_4\gtrsim S_3S_4$. If instead $|A_2|\gtrsim S_3S_4$ then
\[
R_1R_2R_3R_4\ \gtrsim\ S_1S_2^{1/2}S_3^{1/2}\,|A_2|\ \gtrsim\ S_1S_2^{1/2}S_3^{1/2}\cdot S_3S_4
= S_1S_2^{1/2}S_3^{3/2}S_4 .
\]

Thus there is an absolute constant $c_0>0$ with
\[
R_3R_4\ge c_0\,S_3S_4\qquad\text{or}\qquad R_1R_2R_3R_4\ge c_0\,S_1S_2^{1/2}S_3^{3/2}S_4 .
\]
Finally fix $k>0$ to be at most the value used above (to defeat \eqref{iso2}(2)) and also $k<c_0$. Then $R_1\le kS_1$ holds by hypothesis, and the displayed dichotomy yields
\[
R_3R_4> k\,S_3S_4\qquad\text{or}\qquad R_1R_2R_3R_4> k\,S_1S_2^{1/2}S_3^{3/2}S_4,
\]
(the strict inequalities follow from $k<c_0$). This is exactly the desired conclusion.
\qed

\begin{thebibliography}{99}
\bibitem{BT} R.\ Bott and L.\ Tu, \emph{Differential Forms in Algebraic Topology}, Springer.
\bibitem{G3} L.\ Guth, \emph{Notes on Gromov's systolic estimate}, Geom.\ Dedicata \textbf{123} (2006), 113--129.
\bibitem{G} L.\ Guth, \emph{Area-expanding embeddings of rectangles}, arXiv:0710.0403.
\bibitem{GH} L.\ Guth, \emph{Isoperimetric inequalities and rational homotopy invariants}, arXiv:0802.3550.
\end{thebibliography}